\section*{第 \ref{ch:b2:deep-time-climate} 章——地质时期的气候}

\begin{solution}{exo:b2:deep-time-climate:1}
(a) $900 + 37100 + 700 + 3 = \qty{38703}{PgC}$，约为大气中 \qty{591}{PgC} 的 65 倍。
(b) $\num{3.5e6}\times0.012011 = \qty{42000}{PgC}$。(c) $591/42000 = \qty{1.4}{\percent}$。
\end{solution}

\begin{solution}{exo:b2:deep-time-climate:2}
(a) 输入 $11.3 + 3 + 12.3 + 8.1 + 2.0 = \qty{36.7}{Tmol.\text{年}^{-1}}$；输出 $26 + 16 =
\qty{42}{Tmol.\text{年}^{-1}}$，二者在各自的不确定度范围内相等。(b) $(11.3 + 3)/36.7 = 0.39$：
远低于一半；输入的大部分是古老的沉积碳。(c) $\num{3.5e6}/42 = \qty{8.3e4}{\text{年}}$。
\end{solution}

\begin{solution}{exo:b2:deep-time-climate:3}
(a) \ce{Mg2SiO4 + 4 CO2 + 2 H2O -> 2 Mg^{2+} + 4 HCO3^- + SiO2}。(b) 结合 \ce{2 Mg^{2+} + 4 HCO3^-
-> 2 MgCO3 + 2 CO2 + 2 H2O}，净反应为 \ce{Mg2SiO4 + 2 CO2 -> 2 MgCO3 + SiO2}：
每摩尔镁橄榄石消耗两摩尔二氧化碳。(c) $\qty{1}{PgC} = 10^{15}/12.011 =
\qty{8.33e13}{mol}$ 的碳，需要 \qty{4.16e13}{mol} 的镁橄榄石，$4.16\times10^{13}
\times\qty{140.7}{g} = \qty{5.86}{Pg}$，接近六十亿吨。
\end{solution}

\begin{solution}{exo:b2:deep-time-climate:4}
赫南特\omterm{def:b2:orbital-forcing:glacials}{冰期}，445 至 \qty{443}{Ma}，冰室；晚古生代最盛\omterm{def:b2:orbital-forcing:glacials}{冰期}，307 至
\qty{299}{Ma}，冰室；土伦阶，94 至 \qty{89}{Ma}，温室（显生宙中最温暖的时期）；
\omterm{def:b2:deep-time-climate:petm}{古新世–始新世极热事件}，55.9 至 \qty{55.7}{Ma}，温室；早始新世气候最适期，53 至 \qty{49}{Ma}，温室；南极冰盖出现，
约 \qty{34}{Ma}，现代冰室开始；北半球冰盖，
\qty{2.7}{Ma}，冰室。
\end{solution}

\begin{solution}{exo:b2:deep-time-climate:5}
(a) $63000/(8.314\times288^2) = \qty{0.0914}{K^{-1}}$，四舍五入得 0.09。(b) $74000/(8.314
\times288^2) = \qty{0.107}{K^{-1}}$。(c) $\ln1.2/0.09 = \qty{2.03}{K}$，$\ln 1.2/0.107 =
\qty{1.70}{K}$。(d) 在稳态下，风化量必须与脱气量相等，这决定了 $e^{k
\Delta T}$ 进而确定 $\Delta T$；灵敏度只决定产生该升温
所需的二氧化碳量。
\end{solution}

\begin{solution}{exo:b2:deep-time-climate:6}
此处 $k_T = 0.0914$。(a) 精确值 $\exp[7577.5\times(1/288 - 1/293)] = 1.567$；近似值
$e^{0.0914\times5} = 1.579$：偏高 \qty{0.8}{\percent}。(b) 精确值 $e^{2.482} = 11.97$；
近似值 $e^{2.741} = 15.50$：偏高 \qty{30}{\percent}。(c) $\ln(\text{比值}) = (E_{\text
a}/R)[x/T_0^2 - x/(T_0(T_0 + x))] = (E_{\text a}/R)\,x^2/(T_0^2(T_0 + x))$。令其等于 $\ln1.1$：
$x^2 = 1.043\,(288 + x)$，故 $x = \qty{17.9}{K}$。升温超过约 \qty{18}{K} 后
指数形式对风化量的高估超过十分之一。
\end{solution}

\begin{solution}{exo:b2:deep-time-climate:7}
(a) $1444/276 = 5.23$，且 $3.93\log_2 5.23 = \qty{9.4}{W.m^{-2}}$。(b) 几乎恰好等于太阳的
\qty{9}{W.m^{-2}}：两种强迫相互抵消。(c) $(1444 - 276)/419.5 =
\qty{2.8}{ppm}$ 每百万年。\qty{3.4}{ppm} 来自对全部
数据拟合的直线，它对每个估计值都加权，而不是只取平滑曲线两端的数值。
\end{solution}

\begin{solution}{exo:b2:deep-time-climate:8}
(a) 从 $3000/20000 = \qty{0.15}{PgC.\text{年}^{-1}}$ 到 $7000/3000 = \qty{2.3}{PgC.\text{年}^{-1}}$，该范围
包含 IPCC 的 0.3 至 1.5。(b) $35\times12.011/44.01 = \qty{9.6}{PgC.\text{年}^{-1}}$，是\omterm{def:b2:deep-time-climate:petm}{古新世–始新世极热事件}
0.3 至 \qty{1.5}{PgC.\text{年}^{-1}} 的 6 至 32 倍。(c) $3.93\log_2(2000/900) = \qty{4.5}{W.m^{-2}}$，
且 $3\log_2 2.22 = \qty{3.5}{\celsius}$：约为 \qty{5}{\celsius} 的七成。其余部分
必须由其他反馈或其他增温来源（如甲烷）补足。
\end{solution}

\begin{solution}{exo:b2:deep-time-climate:9}
(a) 现在 $1.2\,W_0e^{k\Delta T} = V = W_0$，故 $\Delta T = -\ln 1.2/0.09 = \qty{-2.03}{K}$。
(b) $278\times2^{-2.03/3} = \qty{174}{ppm}$。(c) 不会：在任何稳态下，风化量都等于脱气量，
而脱气量并未改变。风化能力的提高被更冷的气候所抵消。
\end{solution}

\begin{solution}{exo:b2:deep-time-climate:10}
(a) $k = 0.119$：$\tau = 591\times0.693/(0.08\times0.119\times3\times0.138) = \qty{104}{ka}$。
(b) $S = 8$：$\tau = \qty{51.5}{ka}$。(c) 两者都使恒温器反应更快，因为 $\tau \propto
1/(kS)$。$E_{\text a}$ 的不确定度约为两倍，$S$ 在这些时间尺度上的不确定度约为
相近的倍数（3 对 5 至 \qty{8}{K}）：两者影响大致相当。
\end{solution}

\begin{solution}{exo:b2:deep-time-climate:11}
(a) $0.08\times5000 = \qty{400}{PgC}$ 留在大气中：$C/C_0 = 991/591 = 1.677$，即
\qty{466}{ppm}；$\Delta T = 3\log_2 1.677 = \qty{2.24}{K}$。(b) $W - W_0 = 0.138\,(e^{0.09\times2.24}
- 1) = \qty{0.0308}{PgC.\text{年}^{-1}}$，风化增加 \qty{22}{\percent}；$5000/0.0308 =
\qty{162}{ka}$。(c) $\tau = \qty{137}{ka}$，$\tau\ln10 = \qty{315}{ka}$。(d) 升温
只随二氧化碳的对数增长，因此较大的过剩量使风化的增幅小于线性定律所假设的，
移除过程开始得更慢。
\end{solution}

\begin{solution}{exo:b2:deep-time-climate:12}
(a) 在 $M^*$ 附近，$\mathrm d\Delta T/\mathrm dM = Sf/(A_0\ln 2\,(1 + fM^*/A_0)) = Sf/(A_0\ln 2\,
C^*/C_0)$，且 $\mathrm dW/\mathrm dM = kW\,\mathrm d\Delta T/\mathrm dM$，此处 $W = V$；
其倒数即为 $\tau'$。(b) $C^*/C_0 = 2^{2.03/3} = 1.597$：$\tau' = 137.4\times1.597/1.2 =
\qty{183}{ka}$。(c) 更慢：在高二氧化碳浓度下，每增加一拍克碳所引起的升温更小，
因为强迫是对数形式的，风化对其的响应因此更弱。
\end{solution}

\begin{solution}{pb:b2:deep-time-climate:1}
\textbf{1.} $1.3 + 1.3 + 4.5 + 4.2 = \qty{11.3}{Tmol.\text{年}^{-1}}$，$\qty{0.136}{PgC.\text{年}^{-1}}$。
\textbf{2.} $11.5\times0.012011 = \qty{0.138}{PgC.\text{年}^{-1}}$。
\textbf{3.} 二者在约 \qty{2}{\percent} 以内一致：硅酸盐风化平衡了火山
脱气，这正是\omterm{def:b2:deep-time-climate:carbon-cycle}{长期碳循环}所要求的。
\textbf{4.} $\num{5.1e18}/0.02897 = \qty{1.76e20}{mol}$ 的空气；$\num{591e15}/12.011 =
\qty{4.92e16}{mol}$ 的碳；$\num{4.92e16}/\num{1.76e20} = \num{2.80e-4}$，即 \qty{280}{ppm}。
\textbf{5.} $591/0.136 = \qty{4350}{\text{年}}$。
\textbf{6.} $1 + 0.08 = 1.08$：仅多 \qty{8}{\percent}。
\textbf{7.} $63000/(8.314\times288^2) = \qty{0.0914}{K^{-1}}$，四舍五入得 0.09。
\textbf{8.} $e^{0.09} = 1.09$：每开尔文多 \qty{9}{\percent}。
\textbf{9.} $\Delta T = S\log_2(1 + fM/A_0)$；其在 0 处的导数为 $Sf/(A_0\ln2) =
3\times0.08/(591\times0.693) = \qty{5.86e-4}{K.PgC^{-1}}$。
\textbf{10.} $W_0k\times\num{5.86e-4} = 0.138\times0.09\times\num{5.86e-4} = \qty{7.28e-6}{\text{年}^{-1}}$。
\textbf{11.} 其倒数为 $A_0\ln2/(fkSW_0)$，即命题中的 $\tau$，因为
线性化方程为 $\mathrm dM/\mathrm dt = -(\mathrm dW/\mathrm dM)\,M$。
\textbf{12.} 四者均是：$\tau$ 与每一个都成反比。
\textbf{13.} $1/\num{7.28e-6} = \qty{1.37e5}{\text{年}}$。
\textbf{14.} $137\times0.09/0.119 = \qty{104}{ka}$。
\textbf{15.} $k_T = 0.107$：$137\times0.09/0.107 = \qty{115}{ka}$。
\textbf{16.} 短了 1.5 至 7 倍：单箱模型及其常数 $f$ 使
恒温器反应过快。
\textbf{17.} $\num{50e6}/\num{1.37e5} = 364$：恒温器有足够时间在每个
阶段内平衡循环，因此这一下降跟随脱气量或风化能力的缓慢变化。
\textbf{18.} $16\times0.137 = \qty{2.2}{ppm}$。
\textbf{19.} $0.08\times5000 = \qty{400}{PgC}$；$278\times991/591 = \qty{466}{ppm}$。
\textbf{20.} $3\log_2 1.677 = \qty{2.24}{K}$。
\textbf{21.} $0.138\,(e^{0.202} - 1) = \qty{0.0308}{PgC.\text{年}^{-1}}$，\qty{22.3}{\percent}。
\textbf{22.} $137\times\ln10 = \qty{315}{ka}$。
\textbf{23.} 风化可能更为敏感（径流、更大的活化能、更大地球系统灵敏度所带来的
缓慢反馈）；陆地和海洋上的有机碳也可能重新生长，
并且比硅酸盐风化更快地从大气中吸收碳。
\textbf{24.} 在该模型中，\omterm{def:b2:deep-time-climate:thermostat}{硅酸盐风化恒温器}移除过量碳，其
e 折时间约为 \textbf{\qty[text-series-to-math=true]{137}{} 千年}。
\end{solution}
